\section*{ВВЕДЕНИЕ}
\addcontentsline{toc}{section}{ВВЕДЕНИЕ}

В условиях стремительного развития информационных технологий и цифровой экономики особую роль приобретают онлайн-платформы, обеспечивающие взаимодействие пользователей в сети Интернет. Электронные сервисы объявлений, маркетплейсы и платформы аренды становятся неотъемлемой частью современной инфраструктуры бизнеса и повседневной жизни. Они позволяют пользователям быстро находить товары и услуги, заключать сделки и обмениваться информацией в режиме реального времени.

Ключевым элементом подобных систем является эффективная коммуникация между пользователями. Наличие встроенного чата, системы уведомлений и инструментов поддержки напрямую влияет на удобство использования платформы, скорость принятия решений и уровень доверия между участниками взаимодействия. Современные веб-технологии позволяют реализовать обмен сообщениями в режиме реального времени, автоматическую отправку уведомлений по электронной почте, а также внедрение модуля генерации договоров, который автоматически формирует юридически значимые документы на основе параметров сделки (сроки, стоимость, стороны, предмет аренды или купли-продажи), снижая риски ошибок и упрощая оформление отношений между пользователями.

Модуль взаимодействия с пользователями является важной подсистемой веб-приложения, так как обеспечивает информационную связь между арендодателем и арендатором (продавцом и покупателем), информирует о событиях системы и поддерживает непрерывность коммуникации. От корректности его проектирования и реализации зависит стабильность работы всей платформы, масштабируемость системы и удобство конечных пользователей.

Современные компании, разрабатывающие цифровые сервисы, стремятся применять микросервисную архитектуру, позволяющую создавать независимые модули, упрощающие сопровождение и дальнейшее развитие проекта. Выделение сервиса чатов, уведомлений и генерации договоров в отдельный компонент обеспечивает гибкость системы, возможность масштабирования и повышение отказоустойчивости.

\emph{Цель настоящей работы} – разработка модуля взаимодействия с пользователями онлайн-площадки для сдачи вещей в аренду, включающего подсистему уведомлений, подсистему пользовательских чатов и подсистему автоматической генерации договоров аренды. Для достижения поставленной цели необходимо решить \emph{следующие задачи:}
\begin{itemize}
\item провести анализ предметной области онлайн-аренды и определить место модуля коммуникации в общей структуре платформы;
\item выполнить обзор существующих решений в части встроенных мессенджеров, систем уведомлений и автоматической генерации договоров аренды;
\item сформулировать функциональные требования к подсистемам чатов, уведомлений и договоров;
\item спроектировать архитектуру модуля взаимодействия с использованием клиент-серверной модели, REST API и WebSocket-соединений;
\item реализовать подсистему чатов с возможностью обмена текстовыми сообщениями и изображениями в реальном времени;
\item реализовать подсистему уведомлений с доставкой через WebSocket и email-каналы;
\item реализовать подсистему автоматической генерации договоров аренды в формате PDF с демонстрационным механизмом подписания;
\end{itemize}

\emph{Структура и объем работы.} Отчет состоит из введения, 4 разделов основной части, заключения, списка использованных источников, 2 приложений. Текст выпускной квалификационной работы равен \formbytotal{lastpage}{страниц}{е}{ам}{ам}.

\emph{Во введении} сформулирована цель работы, поставлены задачи разработки, описана структура работы, приведено краткое содержание каждого из разделов.

\emph{В первом разделе} рассматривается предметная область онлайн-аренды вещей, проводится классификация онлайн-площадок, анализ существующих решений и аналогов, а также описываются компоненты программно-информационной системы, включая модуль генерации договоров и модуль коммуникации.

\emph{Во втором разделе} описываются функциональные требования к веб-приложению, требования пользователя к интерфейсу и моделируются варианты использования, среди которых формирование и подписание договора аренды, чат между пользователями и управление бронированиями.

\emph{В третьем разделе} на стадии технического проектирования представлены проектные решения, описание разработанных модулей, серверная структура, маршруты API.

\emph{В четвертом разделе} описываются разработанные модули, приведены модульные и интеграционные результаты тестирования разработанных модулей, проверяется корректность разработанных API, оценивается производительность.

В заключении излагаются основные результаты работы, полученные в ходе разработки и делаются выводы о соответствии готового модуля поставленным задачам.

В приложении А представлен графический материал. В приложении Б представлены фрагменты исходного кода. 
